\chapter{Calcitonin}

\section*{What Is This Test?}
The calcitonin test measures the concentration of calcitonin in the blood. Calcitonin is a 32-amino-acid polypeptide hormone synthesized and secreted by the parafollicular C cells (clear cells) of the thyroid gland. Its principal physiologic effects are lowering of serum calcium and phosphate by inhibiting osteoclastic bone resorption and increasing renal calcium and phosphate excretion. It is a counter-regulatory hormone to PTH and 1,25-dihydroxyvitamin D, although its role in normal calcium homeostasis in adults is modest. Calcitonin is the most sensitive and specific tumor marker for medullary thyroid carcinoma (MTC), a tumor that arises from the parafollicular C cells. Basal and stimulated calcitonin levels are used for the diagnosis, staging, treatment monitoring, and recurrence surveillance of MTC. Calcitonin is also elevated in C-cell hyperplasia (a precursor of MTC), in some patients with MEN2 syndromes, and rarely in other neuroendocrine tumors (NETs), small cell lung cancer, breast cancer, and renal cell carcinoma.

\section*{Alternative Names}
\begin{itemize}
  \item Calcitonin Level
  \item Calcitonin Test
  \item Serum Calcitonin
  \item Basal Calcitonin
  \item Stimulated Calcitonin (Pentagastrin or Calcium Stimulation)
\end{itemize}
\section*{Principle}
Calcitonin is measured by immunoassay, typically a chemiluminescent immunoassay (CLIA) or radioimmunoassay (RIA). The assay measures calcitonin monomer, which is the physiologically active form. Calcitonin is released in response to increased serum calcium and by gastrin and pentagastrin. Pentagastrin stimulation testing is the most sensitive method for detecting C-cell hyperplasia and early medullary thyroid carcinoma: after an IV bolus of pentagastrin (0.5 $\mu$g/kg), serum calcitonin is measured at baseline and at 2, 5, and 10 minutes. In normal individuals, the peak stimulated calcitonin is <100 pg/mL. In C-cell hyperplasia, the peak is usually 100--300 pg/mL; in MTC, the peak is typically >300 pg/mL and often >1000 pg/mL. Basal calcitonin >100 pg/mL with a positive pentagastrin stimulation test is highly suspicious for MTC. Calcitonin doubling time is a powerful prognostic marker in patients with metastatic MTC: a doubling time of <6 months predicts poor survival, while a doubling time >2 years predicts an indolent course. Calcitonin is secreted in response to hypercalcemia and is a physiologic calcium-lowering hormone, though its long-term role in adult calcium homeostasis is minimal.

\section*{History}
Calcitonin was discovered in 1961 by Copp and colleagues, who observed a calcium-lowering factor in dog thyroid perfusates \cite{copp1962}. It was initially called calcitonin (tone of calcium). The C cells were identified as the source, and the hormone was purified and sequenced in the 1960s. Calcitonin was used for many years as a therapeutic agent for Paget disease of bone, hypercalcemia, and postmenopausal osteoporosis. It was later recognized that medullary thyroid carcinoma (MTC), a tumor derived from the C cells, produces calcitonin. Basal and stimulated calcitonin measurement became the standard diagnostic and monitoring tool for MTC. The pentagastrin stimulation test was developed in the 1980s. Prophylactic thyroidectomy guided by RET proto-oncogene testing and stimulated calcitonin transformed the management of familial MTC and MEN2 syndromes \cite{wells2015}. Large studies (including the French GETC cohort) established basal calcitonin >100 pg/mL as a strong predictor of MTC and defined calcitonin doubling-time kinetics as a prognostic marker in metastatic disease \cite{barbet2005}.

\section*{Why This Test Is Ordered}
\begin{itemize}
  \item Evaluation of thyroid nodules: some guidelines recommend measuring calcitonin in all patients with thyroid nodules. A basal calcitonin >100 pg/mL is highly suspicious for medullary thyroid carcinoma (MTC) and triggers total thyroidectomy with central neck dissection. Basal calcitonin 20--100 pg/mL is indeterminate and warrants a pentagastrin or calcium stimulation test
  \item Diagnosis of medullary thyroid carcinoma (MTC) and C-cell hyperplasia: elevated basal calcitonin with a positive stimulation test confirms the diagnosis before surgery
  \item Monitoring after total thyroidectomy for MTC: calcitonin should fall to undetectable (<2--10 pg/mL) after complete resection. Persistently detectable calcitonin suggests residual or metastatic disease
  \item Surveillance for recurrence in MTC: rising calcitonin after an undetectable nadir indicates biochemical recurrence. Calcitonin doubling time guides the aggressiveness of further workup (imaging of neck, chest, abdomen, and bones)
  \item Screening and management of MEN2 syndromes (MEN2A, MEN2B, familial MTC): calcitonin is measured in RET-mutation carriers to guide prophylactic thyroidectomy timing and post-operative monitoring. Calcitonin is also measured in patients with MEN2A for pheochromocytoma screening
\end{itemize}

\section*{Primary vs Secondary Test}
Calcitonin is a primary test for the diagnosis and monitoring of medullary thyroid carcinoma. It is a secondary/supplementary test in the evaluation of thyroid nodules (used when MTC is suspected on ultrasound or cytology). It is not used as a general calcium metabolism test.

\section*{How This Test Is Performed}
A healthcare professional draws blood from a vein. For the basal calcitonin test, a single fasting sample is sufficient. For the pentagastrin stimulation test, pentagastrin (0.5 $\mu$g/kg) is given as an IV bolus, and samples are drawn at baseline (time 0) and at 2, 5, and 10 minutes after the injection. In some centers, a calcium infusion stimulation test (2 mg/kg/min elemental calcium over 5--10 minutes) is used instead. The samples are collected in chilled EDTA (lavender-top) tubes and kept on ice until centrifugation, because calcitonin is unstable at room temperature. The specimen should be processed within 2 hours and stored frozen until analysis.

\section*{What Preparation Is Needed}
Fasting is generally recommended for basal calcitonin. The pentagastrin stimulation test requires fasting and cardiac monitoring during the procedure (pentagastrin can cause flushing, nausea, abdominal cramps, and transient hypotension). Patients should stop proton pump inhibitors (PPIs) and H2 blockers (e.g., omeprazole, ranitidine) 3--5 days before testing if possible, because these medications increase gastrin, which stimulates calcitonin release and may elevate baseline levels. Document all current medications.

\section*{Contraindications and Precautions}
\begin{itemize}
  \item Calcitonin is not specific for medullary thyroid carcinoma. It is elevated in C-cell hyperplasia, other neuroendocrine tumors (small cell lung cancer, bronchial carcinoid, pancreatic NET, pheochromocytoma), chronic kidney disease (decreased clearance), hypergastrinemia (pernicious anemia, PPI use, Zollinger-Ellison syndrome), and physiologic states (infants, pregnancy, lactation)
  \item The pentagastrin stimulation test is contraindicated in patients with cardiovascular disease, asthma, severe hypertension, and in pregnancy. It must be performed with IV access and continuous monitoring because of transient flushing, nausea, vomiting, and hypotension
  \item Calcitonin reference ranges are sex-specific: men have higher basal calcitonin than women (approximately 2x). The upper reference limit is typically <10 pg/mL in women and <18 pg/mL in men
  \item A single normal basal calcitonin does not exclude MTC in patients with small tumors (microcarcinoma <1 cm) or in the familial/MEN2 setting, where stimulation testing is more sensitive
  \item Calcitonin assay interference: heterophilic antibodies can cause falsely low or high values; repeat testing in a different assay should be considered if results conflict with the clinical picture
  \item Calcitonin is used to monitor MTC but cannot be used for screening the general population or for diagnosing hypercalcemia
\end{itemize}
\section*{Risks}
Basal testing carries standard venipuncture risks. The pentagastrin stimulation test may cause transient flushing, nausea, abdominal cramping, and lightheadedness (lasting 1--2 minutes). Cardiac monitoring is standard during the test.

\section*{What Happens After the Test}
Results are available within 1--3 days. If basal calcitonin is elevated, the provider correlates it with thyroid ultrasound, cytology, and RET mutation status. If the calcitonin is >100 pg/mL, total thyroidectomy with central neck dissection is typically recommended. After surgery, calcitonin is measured 6--12 weeks post-operatively to confirm a nadir. If calcitonin remains detectable, the patient is evaluated for residual disease. Calcitonin doubling time is calculated from serial measurements and guides the intensity of surveillance.

\section*{Understanding Results}

\subsection*{Below Normal}
\begin{itemize}
  \item \textbf{Undetectable calcitonin (<2 pg/mL) after thyroidectomy:} Complete biochemical cure of MTC. No residual C-cell tissue. The postoperative nadir is the baseline for future surveillance
  \item \textbf{Normal basal calcitonin (sex-specific reference range):} Excludes C-cell disease in most patients with a thyroid nodule, provided the assay is sensitive. In the familial setting, a normal basal level does not exclude early C-cell hyperplasia; stimulation testing is more sensitive
  \item \textbf{Low calcitonin during monitoring:} Stable or falling values indicate no recurrence. Doubling time >2 years predicts an indolent course with favorable prognosis
  \item \textbf{Normal calcitonin with a suspicious thyroid nodule:} MTC is unlikely but not excluded. Cytology (FNA) and ultrasound features guide management
\end{itemize}

\subsection*{Above Normal}
\begin{itemize}
  \item \textbf{Basal calcitonin 20--100 pg/mL:} Indeterminate. Triggers a pentagastrin (or calcium) stimulation test. A peak <100 pg/mL suggests C-cell hyperplasia (premalignant) or assay interference; a peak >100 pg/mL (especially >300 pg/mL) is suggestive of MTC
  \item \textbf{Basal calcitonin >100 pg/mL:} Highly suggestive of MTC (sensitivity ~100\%). Total thyroidectomy with central neck dissection is recommended. The degree of elevation correlates with tumor burden: >100 pg/mL usually corresponds to tumor confined to the thyroid; >500 pg/mL may indicate nodal metastases; >1000 pg/mL may indicate distant metastases
  \item \textbf{Calcitonin >1000 pg/mL (or high stimulated values):} Metastatic MTC should be suspected. Whole-body imaging (CT of neck, chest, abdomen; FDG-PET/CT; and bone scan) is warranted. Calcitonin doubling time is calculated from serial measurements
  \item \textbf{Rising calcitonin in a patient with a goiter or hypergastrinemia:} Consider C-cell hyperplasia, gastrin-mediated stimulation (PPIs, pernicious anemia), or chronic kidney disease. Confirm with a stimulation test and treat the underlying cause; a persistently high calcitonin despite correction suggests C-cell disease
\end{itemize}

\section*{Normal Results}
\begin{itemize}
  \item \textbf{Basal calcitonin (women):} <10 pg/mL (some laboratories <5 pg/mL)
  \item \textbf{Basal calcitonin (men):} <18 pg/mL (some laboratories <15 pg/mL)
  \item \textbf{Peak after pentagastrin stimulation:} <100 pg/mL in normal individuals
  \item \textbf{Post-thyroidectomy (curative surgery for MTC):} undetectable (<2 pg/mL)
\end{itemize}

\section*{Follow-up Tests}
\begin{itemize}
  \item \textbf{Pentagastrin or calcium stimulation test:} For indeterminate basal calcitonin (20--100 pg/mL) or for familial screening. Peak calcitonin >100 pg/mL (especially >300 pg/mL) indicates MTC
  \item \textbf{Thyroid ultrasound and FNA cytology:} For thyroid nodules with elevated calcitonin. FNA with calcitonin measurement in the washout fluid improves diagnostic yield
  \item \textbf{RET proto-oncogene testing (germline):} For all patients with MTC. Germline RET mutations define MEN2A, MEN2B, and familial MTC; prophylactic thyroidectomy timing is guided by the specific mutation risk category (highest risk: MEN2B RET M918T)
  \item \textbf{CEA (carcinoembryonic antigen):} Co-secreted with calcitonin by MTC. CEA is elevated in MTC and is a secondary marker. A rising CEA with stable calcitonin may indicate tumor dedifferentiation and a worse prognosis
  \item \textbf{CT neck/chest/abdomen and FDG-PET/CT:} For staging and restaging when calcitonin suggests metastatic disease (basal >1000 pg/mL or rising levels)
  \item \textbf{Serum calcium, phosphate, and PTH:} To assess the calcium-PTH axis, particularly in MEN2 syndromes where hyperparathyroidism (MEN2A) and pheochromocytoma screening are indicated
\end{itemize}

\section*{References}
\bibliographystyle{unsrtnat}
\bibliography{references}

\clearpage
\section*{Regulatory Status and Authorization Date}
This chapter describes a test category, not a specific commercial product. FDA, CDSCO, EU IVDR, and MHRA approval or registration dates therefore cannot be assigned without the assay manufacturer, product name, instrument, and intended use. For laboratory-developed tests, an individual agency approval date may not exist. See the Preface for the interpretation of regulatory status.

\clearpage
